\begin{afterword}
\summaryhead

The post argues that thinking in terms of ``PR'' makes people and organizations behave worse, while thinking in terms of ``reputation'' or ``honor'' does not. Honor, in the author's sense, means keeping fixed standards and being known to keep them, as the Lannisters are known to pay their debts. PR means predicting what might make people angry and avoiding it.

The author says this second process has no fixed standard, feeds on its own fears, and pushes people toward saying nothing. She also notes that PR seems to need outside experts, while honor is something people can judge for themselves. Her advice: when you catch yourself worrying about PR, ask instead what your reputation or good name requires.

\responsehead

The post has one good observation. People who worry about PR go to outside experts; people who think about their honor judge for themselves. Everything else rests on a distinction the post creates by definition and then treats as a discovery.

``Honor'' is defined as keeping fixed standards and being known to keep them. ``PR'' is defined as predicting what will upset people and avoiding it. With those definitions, honor wins before the argument starts. But in ordinary use the two words name much the same concern: what other people think of you. A company that promises never to sell its users' data, and keeps the promise, is doing PR by any normal meaning of the word. It is also keeping a fixed standard. The gentleman who fought a duel was defending his honor. He was also doing exactly what the post says PR does: acting to head off other people's contempt. The post's own examples cross its own line.

No evidence is offered. The central claims are marked ``I predict'' and ``(I think?)''. No person or organization is described that switched from one way of thinking to the other and did better. The advice is to change a word, which costs nothing and can never visibly fail.

The timing suggests what the post may be for. It appeared on 14 February 2021, the day after the New York Times published Cade Metz's critical article about Slate Star Codex and the rationalist community. The post does not mention the article, so the connection is my inference. If it is right, the post gives a community under press scrutiny a flattering way to describe its own defensiveness: other people do PR; we defend our honor.

In short: a real point about experts, inside a distinction that its own examples dissolve.
\end{afterword}
